\cleardoublepage
\chapter{STUDI LITERATUR}
\label{chap:studi-literatur}

Bab ini membahas konsep yang menjadi dasar pengembangan frontend dashboard WTP
dan membandingkan tugas akhir dengan penelitian terdahulu. 

\section{Pemantauan Instalasi Pengolahan Air}

Instalasi Pengolahan Air atau \textit{Water Treatment Plant} (WTP) menjalankan
serangkaian proses fisika dan kimia yang saling bergantung untuk menurunkan
kandungan partikel tersuspensi, mikroorganisme, dan zat terlarut dalam air
baku. Pada tahap koagulasi, penambahan bahan kimia koagulan seperti Poly
Aluminium Chloride (PAC) menetralkan muatan negatif pada permukaan partikel
koloid sehingga gaya tolak-menolak antarpartikel berkurang dan partikel dapat
saling menempel membentuk flok yang lebih besar. Flok yang terbentuk kemudian
mengendap pada tahap sedimentasi akibat gaya gravitasi, sedangkan partikel
halus yang belum mengendap disaring pada tahap filtrasi. Tahap desinfeksi,
umumnya menggunakan kaporit, bertujuan menonaktifkan mikroorganisme patogen
melalui oksidasi dinding sel maupun komponen metabolik mikroorganisme,
sedangkan penambahan soda ash berfungsi menaikkan nilai pH air untuk mencegah
sifat korosif maupun menjaga efektivitas proses desinfeksi. Karena setiap
tahapan bergantung pada hasil tahapan sebelumnya, perubahan kondisi air baku
pada satu proses dapat memengaruhi efektivitas proses berikutnya sehingga
pemantauan kondisi operasional secara berkelanjutan menjadi bagian penting
dalam menjaga kualitas hasil pengolahan.

Penentuan dosis bahan kimia yang tepat menjadi salah satu aspek kritis dalam
proses tersebut. Dosis yang terlalu rendah menyebabkan proses koagulasi,
desinfeksi, atau koreksi pH tidak berjalan optimal sehingga kualitas air hasil
pengolahan berisiko tidak memenuhi standar yang ditetapkan. Sebaliknya, dosis
yang berlebihan tidak hanya meningkatkan biaya operasional, tetapi juga dapat
menyisakan residu bahan kimia pada air hasil pengolahan. Oleh karena itu,
proses pemberian dosis perlu disertai dengan pemantauan parameter kualitas air
secara berkelanjutan agar penyesuaian dosis dapat dilakukan sesuai dengan
kondisi air baku yang berubah-ubah.

\section{Parameter Kualitas Air}

Parameter kualitas air merupakan indikator yang digunakan untuk menilai kondisi
air selama proses pengolahan. Pada sistem WTP, pemantauan parameter kualitas air
dilakukan secara berkala untuk memastikan bahwa proses pengolahan berjalan
dengan baik dan menghasilkan air yang memenuhi standar yang telah ditetapkan
\autocite{kementerian2023permenkes}. Perubahan nilai parameter dapat menunjukkan adanya perubahan karakteristik air
baku maupun gangguan pada proses pengolahan sehingga informasi tersebut menjadi
dasar bagi operator dalam melakukan evaluasi dan pengambilan keputusan.

Pada tugas akhir ini, \textit{frontend dashboard} menyajikan tiga parameter utama, yaitu
pH, kekeruhan, dan \textit{Total Dissolved Solids} (TDS). Ketiga parameter
tersebut dipilih karena mewakili kondisi kualitas air yang dipantau pada sistem
WTP ITB Jatinangor serta digunakan sebagai dasar evaluasi terhadap hasil
pengolahan.

\subsection{Derajat Keasaman (pH)}

pH merupakan parameter yang menunjukkan tingkat keasaman atau kebasaan suatu
larutan dengan rentang nilai 0 hingga 14. Nilai pH sebesar 7 menunjukkan kondisi
netral, sedangkan nilai di bawah 7 bersifat asam dan nilai di atas 7 bersifat
basa. Dalam proses pengolahan air, nilai pH memengaruhi efektivitas berbagai
tahapan pengolahan, termasuk proses koagulasi dan desinfeksi. Oleh karena itu,
pemantauan pH diperlukan agar operator dapat memastikan proses pengolahan
berlangsung pada rentang yang diharapkan.

Pada \textit{dashboard}, nilai pH ditampilkan dalam bentuk nilai terkini,
grafik tren historis, serta riwayat pengukuran. Penyajian tersebut membantu
operator mengamati perubahan pH dari waktu ke waktu tanpa harus melakukan
pengukuran secara manual.

\subsection{Kekeruhan}

Kekeruhan (\textit{turbidity}) merupakan ukuran banyaknya partikel tersuspensi
yang menyebabkan air tampak keruh. Nilai kekeruhan umumnya dinyatakan dalam
satuan Nephelometric Turbidity Unit (NTU). Semakin tinggi nilai kekeruhan,
semakin banyak partikel yang terdapat di dalam air sehingga proses pengolahan
perlu mendapatkan perhatian lebih lanjut.

Pemantauan kekeruhan penting karena perubahan nilai kekeruhan dapat menunjukkan
perubahan kondisi air baku maupun penurunan efektivitas proses koagulasi,
flokulasi, sedimentasi, atau filtrasi. Pada \textit{dashboard}, data kekeruhan disajikan
dalam bentuk nilai terbaru, grafik perubahan terhadap waktu, dan riwayat
pengukuran sehingga operator dapat memantau kecenderungan perubahan parameter
tersebut.

\subsection{Total Dissolved Solids (TDS)}

\textit{Total Dissolved Solids} (TDS) merupakan parameter yang menunjukkan
jumlah zat padat terlarut di dalam air. Nilai TDS dipengaruhi oleh berbagai
mineral, garam, maupun senyawa lain yang larut dalam air dan umumnya dinyatakan
dalam satuan miligram per liter (mg/L) atau \textit{parts per million} (ppm).

Pemantauan TDS membantu operator memperoleh gambaran mengenai perubahan
kandungan zat terlarut selama proses pengolahan. Meskipun TDS tidak secara
langsung menunjukkan kualitas mikrobiologi air, perubahan nilai TDS dapat
menjadi indikator perubahan karakteristik air baku maupun kondisi proses
pengolahan. Oleh karena itu, parameter ini menjadi salah satu informasi yang
ditampilkan pada dashboard monitoring.

\subsection{Pemanfaatan Parameter pada Dashboard}

Ketiga parameter tersebut ditampilkan secara terintegrasi pada
\textit{dashboard} dalam bentuk kartu ringkasan, grafik tren, analitik sederhana, serta
tabel riwayat pengukuran. Penyajian informasi secara visual membantu operator
membandingkan kondisi antarparameter, mengamati perubahan yang terjadi dalam
periode tertentu, serta mengidentifikasi indikasi kondisi yang memerlukan
perhatian atau tindakan lebih lanjut.

Selain menampilkan nilai hasil pengukuran, \textit{dashboard} juga menerima hasil evaluasi
kepatuhan terhadap batas parameter yang ditentukan oleh \textit{backend}. Dengan demikian,
\textit{dashboard} berperan sebagai media visualisasi informasi, sedangkan proses
penentuan status kepatuhan tetap dilakukan pada sisi \textit{backend} sesuai konfigurasi
yang digunakan sistem.

\section{Internet of Things untuk Pemantauan Air}

\textit{Internet of Things} (IoT) merupakan konsep yang menghubungkan berbagai
perangkat fisik, seperti sensor dan aktuator, ke jaringan komunikasi sehingga
perangkat tersebut dapat mengumpulkan, mengirimkan, dan menerima data secara
otomatis. Pada sistem pemantauan, IoT memungkinkan proses pengambilan data
dilakukan secara berkelanjutan tanpa memerlukan pencatatan manual oleh operator.
Data yang diperoleh dari sensor kemudian dikirim ke sistem pengolahan untuk
disimpan, dianalisis, maupun ditampilkan pada antarmuka pengguna
\autocite{ismail2022iot}.

Dalam sistem pemantauan kualitas air, perangkat IoT umumnya digunakan untuk
mengukur berbagai parameter, seperti pH, kekeruhan, suhu, maupun
\textit{Total Dissolved Solids} (TDS). Sensor melakukan pengukuran secara
berkala, kemudian hasil pengukuran dikirim melalui jaringan komunikasi menuju
\textit{server} atau layanan komputasi yang bertugas mengolah data. Pendekatan tersebut
memungkinkan operator memperoleh informasi kondisi sistem secara lebih cepat
dibandingkan proses pengukuran dan pencatatan manual.

Secara umum, arsitektur sistem pemantauan berbasis IoT terdiri atas empat
lapisan utama, yaitu lapisan akuisisi data (\textit{sensing layer}), lapisan
komunikasi (\textit{communication layer}), lapisan pengolahan data
(\textit{processing layer}), dan lapisan aplikasi (\textit{application layer}),
sebagaimana ditunjukkan pada Gambar~\ref{fig:arsitektur-iot}. Lapisan akuisisi
bertugas memperoleh data dari sensor, lapisan komunikasi mengirimkan data
menuju server, lapisan pengolahan melakukan penyimpanan dan analisis data,
sedangkan lapisan aplikasi menyajikan informasi kepada pengguna melalui
dashboard atau antarmuka lainnya \autocite{ismail2022iot}.

\begin{figure}[H]
  \centering
  \begin{tikzpicture}[
      node distance=0.9cm,
      layer/.style={
        draw, rounded corners, minimum width=8.5cm, minimum height=1.4cm,
        align=center, font=\small
      },
      arr/.style={-{Latex[length=2.5mm]}, thick}
    ]
    \node[layer, fill=blue!8] (sensing) {
      \textbf{Sensing Layer} \\
      Sensor pH, kekeruhan, TDS pada unit ESP32
    };
    \node[layer, fill=blue!15, below=of sensing] (comm) {
      \textbf{Communication Layer} \\
      Transmisi data melalui WiFi menuju server
    };
    \node[layer, fill=blue!22, below=of comm] (proc) {
      \textbf{Processing Layer} \\
      Penyimpanan, pengolahan, dan penyediaan data melalui REST API (\textit{Backend})
    };
    \node[layer, fill=blue!30, below=of proc] (app) {
      \textbf{Application Layer} \\
      Visualisasi data pada \textit{dashboard} operator (\textit{Frontend})
    };

    \draw[arr] (sensing) -- (comm);
    \draw[arr] (comm) -- (proc);
    \draw[arr] (proc) -- (app);
  \end{tikzpicture}
  \caption{Arsitektur Empat Lapisan Sistem Pemantauan Berbasis IoT}
  \label{fig:arsitektur-iot}
\end{figure}

Berbagai penelitian telah memanfaatkan IoT untuk sistem pemantauan kualitas air.
Penelitian oleh \textcite{kumar2024iot} menunjukkan penggunaan sensor pH,
kekeruhan, dan TDS sebagai parameter utama dalam pemantauan kualitas air secara
\textit{real-time}. Data hasil pengukuran dikumpulkan secara otomatis sehingga
perubahan kondisi air dapat diketahui lebih cepat dibandingkan metode
konvensional yang mengandalkan pencatatan manual.

Pada sistem otomasi dosis WTP ITB Jatinangor, frontend tidak berkomunikasi
langsung dengan perangkat sensor maupun protokol komunikasi yang digunakan oleh
perangkat IoT. Seluruh data yang ditampilkan pada dashboard diperoleh dari
backend melalui REST API setelah melalui proses akuisisi, penyimpanan, dan
pengolahan data. Dengan pendekatan tersebut, frontend hanya berperan sebagai
lapisan aplikasi yang bertugas menyajikan informasi kepada operator, sedangkan
seluruh proses komunikasi dengan perangkat IoT menjadi tanggung jawab sistem
backend.

Pemisahan tanggung jawab antara perangkat IoT, \textit{backend}, dan \textit{frontend} memberikan
beberapa keuntungan, seperti kemudahan pengembangan, pemeliharaan sistem, serta
pengurangan keterkaitan antar komponen. \textit{Backend} dapat mengelola proses
pengumpulan data, validasi, dan pengolahan informasi tanpa bergantung pada
implementasi antarmuka, sedangkan \textit{frontend} dapat berfokus pada penyajian data
dan pengalaman pengguna. Arsitektur tersebut juga mempermudah pengembangan
lanjutan apabila di masa mendatang terjadi perubahan perangkat sensor maupun
mekanisme komunikasi yang digunakan.

\section{Platform Penyajian Data Pemantauan}

Penyajian data pemantauan kepada pengguna dapat ditempuh melalui beberapa pendekatan yang telah berkembang dalam
praktik industri maupun penelitian. Subbab ini meninjau cara kerja masing-masing pendekatan sebagai dasar
pertimbangan pada analisis pemilihan solusi di Bab~\ref{chap:analisis-masalah}.

\subsection{Perkakas \textit{Business Intelligence}}
Perkakas \textit{business intelligence} (BI) seperti Power BI dan Tableau bekerja dengan menghubungkan diri ke
sumber data tabular, seperti basis data relasional atau berkas terstruktur, kemudian menyediakan kanvas visual
tempat pengguna menyusun grafik dan laporan melalui konfigurasi tanpa penulisan kode. Data diperbarui melalui
mekanisme penyegaran terjadwal maupun kueri langsung ke sumber data.

Kekuatan pendekatan ini terletak pada kematangan visualisasi, kemudahan penggunaan bagi pengguna non-teknis, dan
kecepatan penyusunan laporan. Keterbatasannya, perkakas BI dirancang untuk analisis data dan bersifat satu arah,
yaitu hanya membaca dan menyajikan data tanpa menyediakan kanal aksi bagi pengguna untuk mengubah keadaan sistem
yang dipantau.

\subsection{Platform IoT Siap Pakai}
Platform IoT siap pakai bekerja dengan menerima data yang dikirimkan perangkat melalui protokol komunikasi
tertentu, menyimpannya pada basis data internal platform, kemudian menyajikannya melalui \textit{dashboard} yang
disusun dari widget bawaan.

ThingSpeak mengorganisasi data ke dalam kanal (\textit{channel}), yaitu wadah data yang setiap bidangnya mewakili
satu parameter. Penyiapannya cepat, tetapi struktur kanal bersifat kaku dan kustomisasi antarmuka sangat terbatas.

Grafana bekerja sebagai lapisan visualisasi yang menghubungkan diri ke sumber data eksternal, umumnya basis data
deret waktu. Grafana tidak menyimpan data secara mandiri, melainkan mengeksekusi kueri ke sumber data setiap kali
panel dirender. Grafana unggul dalam visualisasi telemetri dan memiliki mekanisme \textit{alerting} berbasis
ambang batas yang matang. Keterbatasannya, karena hanya berperan sebagai lapisan visualisasi, seluruh komponen
lain seperti \textit{broker} pesan, basis data, dan manajemen perangkat harus dirakit secara terpisah. Selain itu,
aksi tulis-balik dari antarmuka menuju perangkat tidak tersedia secara bawaan.

ThingsBoard merupakan platform IoT terpadu yang menggabungkan manajemen perangkat, penyimpanan data,
\textit{dashboard}, mesin aturan (\textit{rule engine}), serta manajemen alarm yang dilengkapi mekanisme
pengakuan. Platform ini juga menyediakan widget kendali yang memungkinkan pengguna mengirimkan perintah kembali ke
perangkat. Cakupannya menjadikannya alternatif yang paling mendekati kebutuhan sistem pemantauan industri.
Keterbatasannya, seluruh logika pemrosesan data diarahkan untuk berjalan di dalam mesin aturan platform, sehingga
sistem yang telah memiliki lapisan pemrosesan sendiri perlu menyesuaikan diri dengan model data platform.

\subsection{Pengembangan Antarmuka Web Mandiri}
Pendekatan ketiga adalah membangun antarmuka web secara mandiri yang mengonsumsi layanan \textit{backend} melalui
antarmuka pemrograman aplikasi. Antarmuka bertanggung jawab menarik data, mengelola kesegarannya, dan
menyajikannya, sementara logika pemrosesan tetap berada pada \textit{backend}.

Pendekatan ini menuntut usaha pengembangan yang lebih besar, tetapi memberikan kendali penuh atas alur kerja
pengguna, logika penyajian yang bersifat khusus terhadap domain permasalahan, serta pola integrasi dengan
\textit{backend} yang dikembangkan sendiri. Pendekatan ini umumnya dipilih ketika kebutuhan sistem melampaui
visualisasi telemetri semata, misalnya ketika antarmuka perlu memfasilitasi aksi pengguna yang mengubah keadaan
sistem atau menyajikan evaluasi berbasis regulasi yang bersifat spesifik.

\section{Penelitian Terkait}

Berbagai penelitian telah membahas penerapan teknologi untuk pemantauan kualitas air, mulai dari pengembangan
perangkat \textit{Internet of Things} (IoT), sistem kendali berbasis SCADA, hingga penyajian data melalui
antarmuka web. Subbab ini meninjau cara kerja penelitian-penelitian tersebut sekaligus mengidentifikasi kelebihan
dan keterbatasannya, sebagai dasar penentuan posisi tugas akhir ini.

\subsection{Tinjauan dan Kerangka Kerja Umum}
\textcite{ismail2022iot} menyajikan tinjauan mengenai penerapan teknologi IoT pada sistem pengelolaan air. Kajian
tersebut memetakan arsitektur sistem menjadi lapisan sensor, jaringan komunikasi, platform komputasi, serta
\textit{dashboard} sebagai media visualisasi. Kelebihan penelitian ini terletak pada cakupannya yang luas dan
kerangka arsitektur berlapis yang dapat dijadikan acuan perancangan. Keterbatasannya, pembahasan berhenti pada
tataran arsitektur; lapisan aplikasi hanya disebut sebagai kotak fungsional tanpa penjelasan mengenai bagaimana
antarmuka semestinya dirancang agar benar-benar mendukung kerja operator.

\textcite{alshami2024iot} melakukan tinjauan komprehensif terhadap inovasi IoT dalam pengelolaan air dan air
limbah. Kelebihannya, penelitian ini memetakan perkembangan mutakhir termasuk integrasi kecerdasan buatan pada
sistem pemantauan air. Namun, sebagaimana kajian tinjauan pada umumnya, penelitian ini tidak menghadirkan
implementasi konkret sehingga persoalan praktis pada lapisan antarmuka tidak terungkap.

\subsection{Sistem Pemantauan Kualitas Air Berbasis IoT}
\textcite{kumar2024iot} mengembangkan sistem pemantauan kualitas air berbasis IoT dengan parameter pH, kekeruhan,
dan TDS. Sistem bekerja dengan mengumpulkan pembacaan sensor secara otomatis dan berkala, kemudian mengirimkannya
ke platform pemrosesan untuk ditampilkan. Kelebihannya, penelitian ini membuktikan bahwa ketiga parameter tersebut
memadai sebagai indikator utama kualitas air, yang menjadi dasar pemilihan parameter pada tugas akhir ini.
Keterbatasannya, fokus penelitian berada pada pengembangan perangkat keras dan proses akuisisi data, sedangkan
perancangan antarmuka pengguna tidak dibahas secara rinci.

\textcite{forhad2024wtp} membangun sistem pemantauan kualitas air secara langsung pada instalasi pengolahan air.
Sistem bekerja dengan mengintegrasikan sensor pH, oksigen terlarut, TDS, dan suhu, yang datanya dikirimkan ke
platform terpusat berbasis komputasi awan melalui mekanisme kendali berbasis \textit{Programmable Logic
Controller} (PLC). Kelebihannya, penelitian ini beroperasi pada instalasi nyata, bukan skala laboratorium, serta
telah menyediakan peringatan waktu nyata, pencatatan riwayat, dan pemantauan jarak jauh. Keterbatasannya,
penekanan penelitian berada pada keandalan akuisisi data dan arsitektur perangkat keras; aspek bagaimana informasi
tersebut disajikan kepada operator, termasuk penanganan kondisi ketika data gagal diperbarui, tidak menjadi objek
pembahasan.

\textcite{rahu2023iot} menyusun kerangka kerja yang memadukan IoT dan pembelajaran mesin untuk analisis serta
prediksi kualitas air. Kelebihannya, pendekatan prediktif memungkinkan antisipasi terhadap penyimpangan parameter
sebelum benar-benar terjadi. Keterbatasannya, kompleksitas model prediktif menuntut ketersediaan data historis
yang memadai, dan hasil prediksi tersebut tetap memerlukan antarmuka yang mampu mengomunikasikannya kepada
operator secara dapat ditindaklanjuti, yang tidak menjadi cakupan penelitian tersebut.

\subsection{Antarmuka Operator dan Pemantauan Berbasis Web}
\textcite{low2022scada} menerapkan pemantauan dan pengendalian berbasis web pada sistem distribusi air skala
laboratorium. Sistem bekerja dengan menyajikan status operasional secara terpusat sekaligus memungkinkan
pengendalian jarak jauh melalui peramban. Kelebihannya, penelitian ini membuktikan bahwa antarmuka web mampu
menangani interaksi dua arah, bukan sekadar visualisasi pasif, sehingga menjadi rujukan penting bagi perancangan
kanal aksi operator pada tugas akhir ini. Keterbatasannya, pengujian dilakukan pada skala laboratorium sehingga
belum menghadapi persoalan yang muncul pada instalasi nyata, seperti ketidakstabilan jaringan atau berhentinya
pembaruan data sensor. Penelitian tersebut juga tidak membahas bagaimana antarmuka semestinya berperilaku ketika
pengambilan data gagal, padahal aspek ini menentukan keandalan informasi yang menjadi dasar keputusan operator.

\textcite{azhar2022pwa} menerapkan teknologi \textit{Progressive Web App} (PWA) pada sistem informasi kesehatan
untuk meningkatkan aksesibilitas dan interaksi pengguna. Sistem bekerja dengan memanfaatkan \textit{service
worker} sehingga aplikasi web dapat dipasang pada perangkat dan tetap tanggap meskipun kondisi jaringan tidak
ideal. Kelebihannya, pendekatan ini terbukti memberikan pengalaman setara aplikasi \textit{native} tanpa
memerlukan distribusi melalui toko aplikasi. Keterbatasannya, penerapannya berada pada domain kesehatan sehingga
karakteristik kebutuhan operator industri, seperti pemantauan berkelanjutan dan penanganan kondisi kritis, belum
teruji pada penelitian tersebut.

\subsection{Posisi Tugas Akhir}
Perbandingan antara penelitian terdahulu dengan tugas akhir ini ditunjukkan pada
Tabel~\ref{tab:penelitian-terkait}.

\begin{table}[htbp]
\centering
\caption{Perbandingan penelitian terkait}
\label{tab:penelitian-terkait}
\renewcommand{\arraystretch}{1.25}
\small
\begin{tabular}{|p{1.8cm}|p{2.8cm}|p{3.2cm}|p{3.2cm}|p{3cm}|}
\hline
\textbf{Penelitian} & \textbf{Cara Kerja} & \textbf{Kelebihan} & \textbf{Kekurangan} & \textbf{Celah yang Diisi} \\ \hline

Ismail dkk. (2022) &
Tinjauan arsitektur IoT berlapis untuk pengelolaan air. &
Kerangka arsitektur menyeluruh sebagai acuan perancangan. &
Lapisan aplikasi tidak dibahas secara rinci. &
Merinci perancangan lapisan aplikasi sebagai antarmuka operator. \\ \hline

Alshami dkk. (2024) &
Tinjauan inovasi IoT pada pengelolaan air dan air limbah. &
Memetakan perkembangan mutakhir secara komprehensif. &
Tidak menghadirkan implementasi konkret. &
Menghadirkan implementasi dan verifikasi fungsional. \\ \hline

Kumar dkk. (2024) &
Akuisisi otomatis parameter pH, kekeruhan, dan TDS berbasis IoT. &
Memvalidasi ketiga parameter sebagai indikator utama. &
Berfokus pada perangkat keras; antarmuka tidak dirinci. &
Menyediakan visualisasi, analitik, dan riwayat parameter. \\ \hline

Forhad dkk. (2024) &
Sensor WTP terhubung PLC dan platform awan terpusat. &
Beroperasi pada instalasi nyata dengan riwayat dan peringatan. &
Tidak membahas penyajian data maupun kegagalan pembaruan data. &
Merancang penanganan kondisi data kedaluwarsa. \\ \hline

Rahu dkk. (2023) &
Kerangka IoT dan pembelajaran mesin untuk prediksi kualitas air. &
Memungkinkan antisipasi penyimpangan parameter. &
Hasil prediksi belum dikomunikasikan melalui antarmuka operasional. &
Menyajikan status kepatuhan yang dapat ditindaklanjuti operator. \\ \hline

Low dkk. (2022) &
Pemantauan dan kendali distribusi air berbasis web skala laboratorium. &
Membuktikan kelayakan interaksi dua arah melalui peramban. &
Skala laboratorium; tidak menangani kegagalan pengambilan data. &
Mengembangkan antarmuka pada instalasi nyata dengan penanganan galat. \\ \hline

Azhar dkk. (2022) &
Penerapan PWA pada sistem informasi kesehatan. &
Pengalaman setara aplikasi \textit{native} tanpa toko aplikasi. &
Domain di luar pemantauan industri. &
Mengadopsi PWA untuk pemantauan WTP lintas perangkat. \\ \hline

\end{tabular}
\end{table}

Berdasarkan Tabel~\ref{tab:penelitian-terkait}, terlihat bahwa penelitian terdahulu didominasi oleh pengembangan
perangkat IoT, akuisisi data sensor, maupun kerangka kerja analitik, sedangkan lapisan antarmuka umumnya
diperlakukan sebagai keluaran akhir yang tidak dibahas secara mendalam. Secara khusus, belum ditemukan penelitian
yang membahas bagaimana antarmuka semestinya berperilaku ketika data yang menjadi dasarnya tidak lagi mutakhir,
padahal kondisi tersebut berpotensi menyebabkan operator mengambil keputusan yang keliru.

Tugas akhir ini menempatkan diri pada celah tersebut, yaitu perancangan dan implementasi \textit{frontend
dashboard} yang tidak hanya memvisualisasikan telemetri, tetapi juga menyediakan kanal aksi operator (pengakuan
alarm dan penyesuaian \textit{offset} dosis), menyajikan evaluasi kepatuhan terhadap regulasi nasional, serta
menangani secara eksplisit kondisi ketika data gagal diperoleh atau telah kedaluwarsa.